% !TEX root = ./../../_Thesis.tex

% chapter's came and label
\chapter{Visual Simulation of Refractive Errors} %% on Monochromatic Images}
\label{chap:Methods}

%Visual simulation is a complex, but highly important task. As we discussed, several techniques narrowed the problem by simulating specific factors that affect perception (\eg, glare), while others simplified the problem by simulating only optical aberrations. In 
This chapter describes the approach used for visual simulation of low-order refractive errors. 
Figure \ref{fig:roadmap} illustrates its pipeline, showing equivalent operations specified both in the spatial and in the frequency domain. 
%
Since we are primarily interested in visual acuity, all experiments and discussions presented here are based on monochromatic images. 
%A comprehensive simulation of refractive errors in the context of color vision would require considering color-perception issues which are beyond the scope of this thesis. 
%Additionally, we explain how the technique can be combined with the minimum amount of light perceived by an individual. 
As visual blurring is a depth-dependent phenomenon, we have adopted the simplifying assumption that the observed images are at some constant depth. For this, we used two sets of charts containing standard Sloan letters (Figure~\ref{fig:sloanletters}): black letters on white background, as well as white letters on black background. 
%we need to capture image and depth information of the object. Then, by converting the low-order aberrations of the individual's optical system into a PSF, we can perform a filtering either in the spatial or frequency domain.
%Figure \ref{fig:roadmap} illustrates the pipeline adopted to simulate refractive errors. The equivalent operations are specified both in the spatial and in the frequency domain. 
The following sections provide the details of each step.    

\begin{figure}[h]
	\centering
	\includegraphics[width=1.0\linewidth]{__Images/04/roadmap.png}
	\caption[A roadmap for Fourier optics]{The pipeline for simulating visual aberrations. In the spatial domain one needs the eye's point spread function (PSF), whereas in frequency domain the optical transfer function (OTF) is required \cite{Mollon2003}.} 
	%
	%both in spatial and frequency domain
	%A roadmap for Fourier optics illustrating two paths for computing the retinal image of an object when the point spread function of the eye (PSF), or equivalently, the optical transfer function (OTF), is known \cite{Mollon2003}.}
	\label{fig:roadmap}
\end{figure}

% Capturing Image Information Section
% !TEX root = ./../../_Thesis.tex

\begin{figure}[htb]
	\centering

	\subfigure{
		\includegraphics[width=0.17\linewidth]{__Images/04/C_20-200@4x.png}
	}
	~
	\subfigure{
		\includegraphics[width=0.17\linewidth]{__Images/04/D_20-200@4x.png}
	}
	~
	\subfigure{
		\includegraphics[width=0.17\linewidth]{__Images/04/H_20-200@4x.png}
	}
	~
	\subfigure{
		\includegraphics[width=0.17\linewidth]{__Images/04/K_20-200@4x.png}
	}
	~
	\subfigure{
		\includegraphics[width=0.17\linewidth]{__Images/04/N_20-200@4x.png}
	}
	~
	\subfigure{
		\includegraphics[width=0.17\linewidth]{__Images/04/O_20-200@4x.png}
	}
	~
	\subfigure{
		\includegraphics[width=0.17\linewidth]{__Images/04/R_20-200@4x.png}
	}
	~
	\subfigure{
		\includegraphics[width=0.17\linewidth]{__Images/04/S_20-200@4x.png}
	}
	~
	\subfigure{
		\includegraphics[width=0.17\linewidth]{__Images/04/V_20-200@4x.png}
	}
	~
	\subfigure{
		\includegraphics[width=0.17\linewidth]{__Images/04/Z_20-200@4x.png}
	}
	
	\caption{Standard Sloan Letters.}
	\label{fig:sloanletters}
\end{figure}

% section's Name and Label
\section{Target Images and Capture Setup}
\label{sec:CapturingImageInformation}

We have created images of Sloan letters with values ranging from -0.3 to 1.0 in steps of 0.1 in the LogMAR (Logarithm of the Minimum Angle of Resolution) scale \cite{Bailey1976}. Such an interval corresponds to the range from 20/10 to 20/200, respectively, in the Snellen scale. The LogMAR  scale provides a more accurate estimate of visual acuity when compared to other charts (\eg, Snellen), being the recommended one for research settings. 
Our target images were created according to Equation~\ref{eq:lettersize} for testing vision from three feet away. The individual letters were rendered using the vector graphics capabilities of Inkscape and the Sloan PostScript fonts provided by \citet{Pelli1988} 
%. Using the standard Sloan letters, we generated files at 360 dpi, containing black letters on white background 
(Figure~\ref{fig:sloanletters}).
% as well as white letters on black background. 
At the prescribed distance, the ratio between one pixel and one arc minute is 1:1, that is, the letters with a LogMAR value of 0 (or Snellen fraction 20/20) has exactly 5 pixels of height. For the purpose of our simulations, each black (white) letter (also called an {\it optotype}) was placed against a $113 \times 133$-pixel black (white) square.   
Since $1~degree = 60~arc~minutes$, each such square covers a total field of view (FOV) of 1,88� $\times$ 1,88�. The conversion from {\it Snellen decimal acuity} values  
%and vice-versa are 
to LogMAR values is presented in Equation \ref{eq:dec2logmar}. A Snellen decimal acuity value is the decimal representation
of the equivalent Snellen ratio (\eg, Snellen ratios of 20/20 and 20/40 correspond to Snellen decimal acuity values of 1.0 and 0.5, respectively). 
% and \ref{eq:logmar2dec}, respectively.

\begin{equation}
	\centering
	\label{eq:lettersize}
	letter~size_{mm} = \left \{  \tan\left [ deg2rad \left (\frac{5}{60} \right ) \right ] \times (chart~distance_{mm}) \times (10^{-LogMAR})^{-1} \right \}
\end{equation}

\begin{equation}
	\centering
	\label{eq:dec2logmar}
	LogMAR = -\log_{10} (Snellen~decimal~acuity)
\end{equation}

%\begin{equation}
	%\centering
	%\label{eq:logmar2dec}
	%Snellen~Decimal~Acuity = 10^{-LogMAR}
%\end{equation}


%In order to capture images of real scenes, 
We have prepared white- and black-background LogMAR charts containing Sloan letters specifically designed for a viewing distance of three feet (Figure~\ref{fig:etdrs_tables}). The charts were printed at 360 dpi on white paper using a laser printer. We then took pictures of the charts with a DSLR camera. The camera was placed at three feet (91.44 cm) from the chart, with focal length set to 18mm. Since images acquired using this setup respect the 1:1 ratio between pixels and arc minutes, one can crop the squares containing the individual optotypes for further processing.

%\begin{lstlisting}[caption = {For educational purposes}]
%function [ L, P, pt, D ] = calc_Sloan( distance_ft, va_snellen )
%	distance_mm   = ft2mm(distance_ft);
%	syms x;
%	eqn = tan(deg2rad(5/60)) == x / (distance_mm * va_snellen);
%	tmp = solve(eqn, x);
%	x   = double(vpa(tmp));
%	
%	L  = x;					% letter size in mm
%	P  = x/5;				% MAR    size in mm
%	pt = mm2psPT(L);		% font   size (Pages, MacOS editor)
%	D  = distance_mm / 10;  % /10   to return distance in cm
%end
%\end{lstlisting}


\begin{figure}[htb]
	\centering

	\subfigure[]{
		\includegraphics[width=0.45\linewidth]{__Images/04/ETDRS_Chart_W.png}
	}
	~
	\subfigure[]{
		\includegraphics[width=0.45\linewidth]{__Images/04/ETDRS_Chart_B.png}
	}
	
	\caption[LogMAR Charts]{LogMAR charts printed at 360 dpi on white paper using a laser printer. (a) Black letters on white background and (b) white letters on black background. The top row of each table corresponds to Snellen 20/200 (LogMAR +1.0) visual acuity when viewed from three feet. The bottom row corresponds to Snellen 20/10 (LogMAR -0.3) visual acuity when viewed from three feet away.}
	\label{fig:etdrs_tables}
\end{figure}

% Specifying Aberrations Section
\input{__Text/04_Methods/420.tex}

% Image Filtering Section
% !TEX root = ./../../_Thesis.tex

% section's Name and Label
\section{Image Filtering}
\label{sec:Image Filtering}

Given $S$, $C$, $R$, and $\phi$,  
one can obtain the effective aberration function as $kW_{(x,y)}$, where $k$ is the spherical wavenumber (\ie, $k = 2\pi/\lambda$), and 
$W_{(x,y)}$ is the wavefront aberration function expressed using the Zernike polymials. For the case of low-order aberrations, $W_{(x,y)}$ is defined by Equation~\ref{eq:W}, which takes into account oblique astigmatism, defocus, and vertical astigmatism. 
%\begin{equation}
%W_{(x,y)} = \sum_{i=-1}^1 c_{2}^{2i} \, Z_{2}^{2i}_{(x,y)}.
%\end{equation}
 $\lambda = 550nm$ is a standard wavelength used for monochromatic simulation \cite{Dai2008}.
 The pupil function $P_{(x,y)}$ is a binary  function that evaluates to 1 inside the projected aperture, and 0 outside it. According to \citet{Goodman2005}, the {\it generalized pupil function} $\mathbb{P}_{(x,y)}$ is given by:
\begin{equation}
	\centering
	\label{eq:generalizedpupilf}
	\mathbb{P}_{(x,y)} = P_{(x,y)} \exp[j*k*W_{(x,y)}],
\end{equation}
%
where $j = \sqrt{-1}$. Note that $\mathbb{P}_{(x,y)}$ is a complex number. One can obtain the point spread function of the optical system as the power spectrum of $\mathbb{P}$, \ie,  $PSF = \left | \mathcal{F}(\mathbb{P}) \right | ^ {2}$, where $\mathcal{F}$ is the Fourier transform operator.  Given the PSF and 
an input image $I$, one can simulate the view of $I$ through the given optical system computing the 2-D convolution
$O = PSF \otimes I$.
% and compute the two-dimensional convolution of $I$ and $\overline{PSF}$ in order to obtain the output image $O$. This is the simpler path for computing retinal images of a Sloan letter $S$ for a fixed depth, carried out entirely in the spatial domain. 
A more efficient computation of $O$ can be obtained in the frequency domain (this is illustrated by purple arrows in Figure~\ref{fig:roadmap}). In that case, $O = \mathcal{F}^{-1}(\mathcal{F}(I) * OTF)$, where $OTF = \mathcal{F}(PSF)$ is the {\it the optical transfer function} and $*$ is the element-wise multiplication. 
%
%the product of the object spectrum $\mathcal{F}(I)$ and the optical transfer function $OTF = \mathcal{F}(PSF)$ is the image spectrum, from which the image itself can be obtained by inverse Fourier transform $\mathcal{F}^{-1}(I * \overline{PSF}).

% Validation Section
% !TEX root = ./../../_Thesis.tex

% section's Name and Label
\section{Validation}} 
\label{sec:Validation}

To validate the visual simulation results of the refractive errors, we use a DSLR camera (Canon model EOS Rebel T3 with an 18-55mm zoom lens). The camera represents a perfect eye (\ie, without refractive aberrations). We place additional lenses in front of the camera's optical system to induce low-order aberrations (\ie, myopia, hyperopia, and astigmatism). Such lenses are placed on a support fixed to a UV filter attached to the main lens. Figure~\ref{fig:extralens} shows the camera with an additional +1.0 diopter lens attached to it. The support can hold up to three additional lenses. 
%
%results from the image filtering process, we've used additional lenses to induce low-order aberrations (\ie, myopia, hyperopia, and astigmatism) into the DSLR camera's optical system (Figure~\ref{fig:extralens}). 

For our simulations, we use a simplified eye model adjusted to the camera's settings to achieve consistent results between them.
% with the ones captured by the camera. 
More specifically, we make sure that the \emph{f-number} (\ie, the ratio of the camera lens' focal length $f$ to the diameter $D$ of its aperture):
\begin{equation}
	\centering
	\label{eq:fnumber}
	f_{number} = \frac{f}{D}
\end{equation}
\noindent
is the same for the camera and the eye model. For the experiments shown in the thesis, we fixed the focal length of the camera's main lens to 18mm (regardless of the use of additional lenses).
% to captured images (with or without extra lenses) with 18mm focal length and 
Thus, for instance, given \emph{f-number} values of $4.0$, $4.5$ and $5.0$, the corresponding camera lens aperture values are
%entrance pupil for each possible combination is
 4.5mm, 4.0mm and 3.6mm, respectively. 
Our simplified eye model (Figure~\ref{fig:syntheticeye}) has an axial diameter of 18mm. The crystalline lens causes the nodal point {\bf N} to be behind the crystalline. Thus, the eye model's effective focal length is 13.5mm: $f_{eye} = 18mm / \eta_{eye} =  18mm / 1.333 = 13.5 mm$, where $\eta_{eye}$ is the index of refraction of the eye. As a result, the eye model's pupil size (equivalent of the camera's lens aperture) needs to be rescaled to maintain the same \emph{f-number} value as the camera. Table~\ref{table:pupildiameter} shows the corresponding values of the equivalent camera apertures and pupil diameters. The simulation results shown in Chapter~\ref{chap:Results} were obtained for $f/5.0$ (third row of Table~\ref{table:pupildiameter}), although other values could have been used. 

%we must rescale the entrance pupil's diameter (Table~\ref{table:pupildiameter}) in order to maintain the same \emph{f-number} in simulations of retinal images. 

\begin{table}[h]
	\centering
	\caption{Camera apertures and pupil diameters for various f-numbers.}
	\label{table:pupildiameter}
	\begin{tabular}{rcc}
		\multicolumn{1}{l}{\bf f-number} & {\bf DSLR Camera (18mm focal length)} & {\bf Synthetic Eye (13.5mm focal length)} \\
		{\bf }               & {\bf aperture}     & {\bf pupil diameter}     \\
		{\bf $f/4.0$}         & 4.5mm                    & 3.4mm                        \\
		{\bf $f/4.5$}         & 4.0mm                    & 3.0mm                        \\
		{\bf $f/5.0$}         & 3.6mm                    & 2.7mm                        
	\end{tabular}
\end{table}

\begin{figure}[h]
	\centering
	
	\subfigure[]{
		\includegraphics[width=0.55\linewidth]{__Images/04/camera.jpg}
		\label{fig:extralens}
	}
	~
	\subfigure[]{
		\includegraphics[width=0.35\linewidth]{__Images/04/eyemodel.png}
		\label{fig:syntheticeye}
	}
	
	\caption[Optical systems used in the validation process]{Optical systems used in the validation process: (a) Canon EOS Rebel T3 with apparatus to add up to three extra lenses. Focal lens set to 18mm. (b) Simplified eye model with effective focal length of 13.5mm. {\bf N} is the nodal point.}
	\label{fig:camera}
\end{figure}

% ABCD Matrix Subsection
% !TEX root = ./../../_Thesis.tex

% section's Name and Label
\subsection{Vertex Distance and Ray Transfer Matrix}
\label{subsec:ABCDMatrix}

\begin{figure}[!b]
	\centering

	\subfigure[]{
		\includegraphics[width=0.40\linewidth]{__Images/04/Zeiss_BackVD.png}
		\label{fig:vertexdistance}
	}
	~
	\subfigure[]{
				\includegraphics[width=0.55\linewidth]{__Images/04/RTM.png}
		\label{fig:rtm_scheme}
	}
	
	\caption[Vertex distance and our optical setup]{Vertex distance. (a) Typical eyeglasses vertex distance of 12mm. (b) Our camera setup with a distance of 10mm between DSLR's main lens and the additional one.}
	\label{fig:rtm}
\end{figure}

The optical power of a lens prescribed for correcting low-order aberrations varies according to the distance from the lens to the cornea, also known as {\it vertex distance} (Figure~\ref{fig:vertexdistance}). 
%While the vertex distance from the cornea to a contact lens is 0mm, its distance to an eyeglass is usually 12mm. 
To compensate for the spacing between camera's main lens and the additional ones (Figure~\ref{fig:rtm_scheme}), we use a {\it ray transfer matrix} (RTM) formulation 
%to describe the light path through the optical system 
\cite{Glytsis2014}.
%provides a step-by-step description of the process involved in computing a ray transfer matrix.
%This calculus, demonstrated step by step by \citet{Glytsis2014}, implicates in the determination of a new focal length and magnification value for the focal system consisting of two lens (\review{see Appendix XYZ}).
%Annex~\ref{chap:Annex1} discusses how to obtain the new focal length and magnification value for an optical system consisting of multiple lenses.
The RTM representing two thin lenses separated by a distance $d$ can be obtained multiplying three matrices:
%eq:W
 a thin lens matrix (that approximates the DSLR's optical system by a single thin lens), a distance $d$ propagation matrix, and a thin lens matrix representing our additional lenses:
%
\begin{equation}
	\centering
	\label{eq:RTM_tl}
	\begin{array} {lcl} \begin{bmatrix} A_{TL} & B_{TL}\\ C_{TL} & D_{TL} \end{bmatrix} 
	& = & \begin{bmatrix} 1 & 0\\ -\frac{1}{f_{camera}} & 1 \end{bmatrix}
	\begin{bmatrix} 1 & d\\ 0 & 1 \end{bmatrix}
	\begin{bmatrix} 1 & 0\\ -\frac{1}{f_{lens}} & 1 \end{bmatrix} \\ \\ 
	& = & \begin{bmatrix}
	1 - \dfrac{d}{f_{camera}} & d\\
	\left (  \dfrac{\dfrac{d}{f_{lens}}-1}{f_{camera}} - \dfrac{1}{f_{lens}}\right ) & 1 - \dfrac{d}{f_{lens}}
	\end{bmatrix}. 
	\end{array}
\end{equation}
%
Here $f_{camera}$ is the DLSR camera focal length (\ie, 18mm in our case), and $f_{lens}$ is the focal length of the (combined set of) additional lens(es).
The image captured by the resulting optical system is formed at a distance $x$ behind the DSLR camera's optical system. Assuming we want to capture the image of an infinitely far away object (\eg, at distance $s = 10^{20}$mm from the camera), the overall RTM can computed as:
%
\begin{equation}
	\centering
	\label{eq:RTM_final}
	\begin{bmatrix} A & B\\ C & D \end{bmatrix}
	=
	\begin{bmatrix} 1 & x\\ 0 & 1 \end{bmatrix}
	\begin{bmatrix} A_{TL} & B_{TL}\\ C_{TL} & D_{TL} \end{bmatrix}
	\begin{bmatrix} 1 & s\\ 0 & 1 \end{bmatrix}.
\end{equation}
%
Since a set of parallel rays (of an infinitely far away object) are focused by a lens to its focal point, one concludes that $x$ should indeed be the focal length $f_{cam+lens}$ of the compounded optical system comprised by the camera's main lens plus the additional one. 
%converge to the optical system's focal point, 
By letting $B = 0$, 
%\review{(justificar por que isso)}, 
one can solve for $x$, obtaining:
%To achieve optical system's new focal length, the ABCD matrix element $B$ must become zero. When $B = 0$, the value of $x$ gives the new focal length and the value of $A$ gives the overall magnification $m$ of the image (\ie, $m = A(x)$).
%
\begin{equation}
x = f_{cam+lens} = \dfrac{(d  + s) \times (f_{camera} \times f_{lens}) - (d \times f_{lens} \times s) }
{(d - f_{lens}) \times f_{camera} + (f_{camera} + f_{lens} - d )\times s}. 
\label{RTMx}
\end{equation}
%
%\begin{subequations}
 %\begin{align}
  %B = 0 \label{RTMb} \\
  %d + s\left (  x\left (\dfrac{ \dfrac{d}{f_{lens}}-1}{f_{camera}} - \dfrac{1}{f_{lens}}\right )-\dfrac{d}{f_{camera}}+1\right ) -x\left (\dfrac{d}{f_{lens}}-1\right ) = 0 \label{RTMb0} \\
  %x = \dfrac{d~f_{camera}~f_{lens} - d~f_{lens}~s+f_{camera}~f_{lens}~s}{d~f_{camera}-f_{camera}~f_{lens}-d~s+f_{camera}~s+f_{lens}~s} \label{RTMx}
 %\end{align}
%\end{subequations}
%
Since 1 diopter = 1/meter, and $f_{cam+lens}$ is expressed in mm, the dioptric power of the resulting compounding optical system is given by:
%The new optical system's dioptric power is given by:
\begin{equation}
	\centering
	\label{eq:newD}
	diopt_{cam+lens}= \frac{1}{f_{cam+lens} \times 10^{-3}} =  \frac{10^3}{f_{cam+lens}} D.
\end{equation}
%
Table~\ref{table:newDioptricPower} shows the actual increase in dioptric power that result from placing additional lenses with different powers in front of the camera's main lens, considering a vertex distance of $10mm$. Thus, for instance, when placing a +1.0D lens in front of the camera's main lens, we are in fact inducing myopia of 1.0101D. Therefore, in order to obtain an image comparable to the one captured by the camera, our simulation should compute a wavefront aberration corresponding to 1.0101D of myopia.    
%
\begin{table}[!h]
\centering
\caption{Actual increase in dioptric power obtained by placing additional lenses with various powers in front of the camera's main lens considering a vertex distance of 10mm.}
\label{table:newDioptricPower}
\begin{tabular}{cc}
\hline
{\bf Additional Lens' dioptric power} & {\bf Actual dioptric power} \\ \hline
0.0000 D                              & 0.0000 D                              \\
1.0000 D                              & 1.0101 D                              \\
2.0000 D                              & 2.0408 D                              \\
3.0000 D                              & 3.0928 D                              \\
4.0000 D                              & 4.1667 D                              \\ \hline
\end{tabular}
\end{table}

% Optical Simulation Section
% !TEX root = ./../../_Thesis.tex

% section's Name and Label
\subsection{Comparison of Simulated Results with Ground Truth}
\label{sec:OpticalSimulation(Results)}

This section compares our simulated results with an optical ground truth, obtained by capturing images of the LogMAR charts shown in Figure~\ref{fig:etdrs_tables}.
%vision chart containing Sloan letters produced using Equation \ref{eq:lettersize} and placed at three feet from the camera. 
Whenever we reference to the dioptric power of additional lenses, our simulations account for the values described in Table~\ref{table:newDioptricPower}.  
%And even though the final arbiter of image quality is the human viewer, 
To objectively evaluate the quality of the simulated results, we use three objective metrics: the Structural Similarity Image Metric (SSIM)~\cite{Wang2004}, the Peak Signal-to-Noise Ratio (PSNR), and the Absolute Difference (AD) of the pixelwise differences between the captured and simulated images.
The SSIM metric measures image degradation perceived as change in structural information. 
%It produces a number between -1.0 (poor similarity) and 1.0 (high similarity). 
It is calculated for each pixel of a given image with respect to some reference image, based on its relationship to other pixels in an 11-by-11 neighborhood. PSNR is a popular metric in image processing for assessing the quality of image reconstruction and compression. It is often expressed using a decibel scale, and computed as  
%
\begin{equation}
	\centering
	\label{eq:PSNR}
	PSNR = 10\log_{10}\left (\frac{peakval^2}{MSE}\right ),
\end{equation}
where
\begin{equation}
	\centering
	\label{eq:MSE}
	MSE = \frac{1}{mn}\sum^{m-1}_{i=0}\sum^{n-1}_{j=0}(I_{ref}(i,j) - I(i,j))^2,
\end{equation}
and $I$ is an image being compared to a reference image $I_{ref}$, both with the same dimensions $m \times n$. 
$peakval$ is the maximum possible value for a pixel. For instance, for a grayscale image using 8-bits per pixel, $peakval=255$. 
   

%
The optical simulation described in this thesis was implemented using MATLAB Student Version (R2014a).
% environment has been chosen to implement the simulation technique described in this thesis. 
%
%As described in Section \ref{sec:Validation}, we've used additional lenses to induce low-order aberrations (\ie, myopia, hyperopia, and astigmatism) into a DSLR camera's optical system. This allow us to 
Figures~\ref{fig:comparison_myopic_wb} and \ref{fig:comparison_myopic_bw} compare images of a letter from the LogMAR charts with 
%Sloan letters against a 
white and black background, respectively, captured by the DSLR camera (top row) against the results of our simulations (second row). 
%of myopic views of the chart. 
The images in the top rows were captured by the camera with extra lenses, ranging from 0 to +4 diopters, in steps of 1 diopter. The second rows show the images produced using our simulation and considering the adjustments in dioptric power required to account for the 10mm spacing between the camera's main lens and the additional one (Table~\ref{table:newDioptricPower}). Our simulations were applied to the image captured by the camera without any extra lens (\ie, camera +0.00 D). The third and fourth rows of these figures show visual representations of the SSIM an AD metrics, respectively. 
%Also the comparison when using the SSIM metric (third row) or AD metric (bottom row).

% Myopic 		- Figures
% !TEX root = ./../../../_Thesis.tex

\begin{figure*}[!t]
	\centering

	\begin{tabular}{@{}r@{ } c@{ } c@{ } c@{ } c@{ } c }
	&
	\small{-0.00 D} &
	\small{-1.00 D} &
	\small{-2.00 D} &
	\small{-3.00 D} &
	\small{-4.00 D} & \\

	\begin{sideways} \parbox[b]{20mm} {Camera} \end{sideways} &
	\includegraphics[width=0.185\textwidth]{__Images/05/BW_N(20-200)_-0D_to_-4D/bw_N_20-200_Camera-0,00D(lens).png} &
	\includegraphics[width=0.185\textwidth]{__Images/05/BW_N(20-200)_-0D_to_-4D/bw_N_20-200_Camera-1,00D(lens).png} &
	\includegraphics[width=0.185\textwidth]{__Images/05/BW_N(20-200)_-0D_to_-4D/bw_N_20-200_Camera-2,00D(lens).png} &
	\includegraphics[width=0.185\textwidth]{__Images/05/BW_N(20-200)_-0D_to_-4D/bw_N_20-200_Camera-3,00D(lens).png} &
	\includegraphics[width=0.185\textwidth]{__Images/05/BW_N(20-200)_-0D_to_-4D/bw_N_20-200_Camera-4,00D(lens).png} \\

	\begin{sideways} \parbox[b]{20mm} {Simulation} \end{sideways} &
	\includegraphics[width=0.185\textwidth]{__Images/05/BW_N(20-200)_-0D_to_-4D/bw_N_20-200_Camera-0,00D(simulated).png} &
	\includegraphics[width=0.185\textwidth]{__Images/05/BW_N(20-200)_-0D_to_-4D/bw_N_20-200_Camera-1,00D(simulated).png} &
	\includegraphics[width=0.185\textwidth]{__Images/05/BW_N(20-200)_-0D_to_-4D/bw_N_20-200_Camera-2,00D(simulated).png} &
	\includegraphics[width=0.185\textwidth]{__Images/05/BW_N(20-200)_-0D_to_-4D/bw_N_20-200_Camera-3,00D(simulated).png} &
	\includegraphics[width=0.185\textwidth]{__Images/05/BW_N(20-200)_-0D_to_-4D/bw_N_20-200_Camera-4,00D(simulated).png} \\

	\begin{sideways} \parbox[b]{20mm} {Local~SSIM} \end{sideways} &
	\includegraphics[width=0.185\textwidth]{__Images/05/BW_N(20-200)_-0D_to_-4D/bw_N_20-200_Camera-0,00D(comparison).png} &
	\includegraphics[width=0.185\textwidth]{__Images/05/BW_N(20-200)_-0D_to_-4D/bw_N_20-200_Camera-1,00D(comparison).png} &
	\includegraphics[width=0.185\textwidth]{__Images/05/BW_N(20-200)_-0D_to_-4D/bw_N_20-200_Camera-2,00D(comparison).png} &
	\includegraphics[width=0.185\textwidth]{__Images/05/BW_N(20-200)_-0D_to_-4D/bw_N_20-200_Camera-3,00D(comparison).png} &
	\includegraphics[width=0.185\textwidth]{__Images/05/BW_N(20-200)_-0D_to_-4D/bw_N_20-200_Camera-4,00D(comparison).png} \\
	
	\begin{sideways} \parbox[b]{20mm} {Difference} \end{sideways} &
	\includegraphics[width=0.185\textwidth]{__Images/05/BW_N(20-200)_-0D_to_-4D/bw_N_20-200_Camera-0,00D(diff).png} &
	\includegraphics[width=0.185\textwidth]{__Images/05/BW_N(20-200)_-0D_to_-4D/bw_N_20-200_Camera-1,00D(diff).png} &
	\includegraphics[width=0.185\textwidth]{__Images/05/BW_N(20-200)_-0D_to_-4D/bw_N_20-200_Camera-2,00D(diff).png} &
	\includegraphics[width=0.185\textwidth]{__Images/05/BW_N(20-200)_-0D_to_-4D/bw_N_20-200_Camera-3,00D(diff).png} &
	\includegraphics[width=0.185\textwidth]{__Images/05/BW_N(20-200)_-0D_to_-4D/bw_N_20-200_Camera-4,00D(diff).png} \\

	\end{tabular}
	
	\caption[Comparisons of our simulated results against ground truth obtained with a hyperopic camera]{Comparisons of our simulated results against ground truth obtained with a hyperopic camera. These images correspond to a Snellen ratio of 20/200. (top row) Images captured using the DSLR camera with extra lenses varying from 0.0 to 4.0 diopters. (second row) Our simulated results. (third row) SSIM metric results. (fourth row) AD metric.}
	\label{fig:comparison_hyperopic_bw}
\end{figure*}
	% fig:comparison_myopic_wb
% !TEX root = ./../../../_Thesis.tex

\begin{figure*}[!t]
	\centering

	\begin{tabular}{@{}r@{ } c@{ } c@{ } c@{ } c@{ } c }
	&
	\small{-0.00 D} &
	\small{-1.00 D} &
	\small{-2.00 D} &
	\small{-3.00 D} &
	\small{-4.00 D} & \\

	\begin{sideways} \parbox[b]{20mm} {Camera} \end{sideways} &
	\includegraphics[width=0.185\textwidth]{__Images/05/BW_N(20-200)_-0D_to_-4D/bw_N_20-200_Camera-0,00D(lens).png} &
	\includegraphics[width=0.185\textwidth]{__Images/05/BW_N(20-200)_-0D_to_-4D/bw_N_20-200_Camera-1,00D(lens).png} &
	\includegraphics[width=0.185\textwidth]{__Images/05/BW_N(20-200)_-0D_to_-4D/bw_N_20-200_Camera-2,00D(lens).png} &
	\includegraphics[width=0.185\textwidth]{__Images/05/BW_N(20-200)_-0D_to_-4D/bw_N_20-200_Camera-3,00D(lens).png} &
	\includegraphics[width=0.185\textwidth]{__Images/05/BW_N(20-200)_-0D_to_-4D/bw_N_20-200_Camera-4,00D(lens).png} \\

	\begin{sideways} \parbox[b]{20mm} {Simulation} \end{sideways} &
	\includegraphics[width=0.185\textwidth]{__Images/05/BW_N(20-200)_-0D_to_-4D/bw_N_20-200_Camera-0,00D(simulated).png} &
	\includegraphics[width=0.185\textwidth]{__Images/05/BW_N(20-200)_-0D_to_-4D/bw_N_20-200_Camera-1,00D(simulated).png} &
	\includegraphics[width=0.185\textwidth]{__Images/05/BW_N(20-200)_-0D_to_-4D/bw_N_20-200_Camera-2,00D(simulated).png} &
	\includegraphics[width=0.185\textwidth]{__Images/05/BW_N(20-200)_-0D_to_-4D/bw_N_20-200_Camera-3,00D(simulated).png} &
	\includegraphics[width=0.185\textwidth]{__Images/05/BW_N(20-200)_-0D_to_-4D/bw_N_20-200_Camera-4,00D(simulated).png} \\

	\begin{sideways} \parbox[b]{20mm} {Local~SSIM} \end{sideways} &
	\includegraphics[width=0.185\textwidth]{__Images/05/BW_N(20-200)_-0D_to_-4D/bw_N_20-200_Camera-0,00D(comparison).png} &
	\includegraphics[width=0.185\textwidth]{__Images/05/BW_N(20-200)_-0D_to_-4D/bw_N_20-200_Camera-1,00D(comparison).png} &
	\includegraphics[width=0.185\textwidth]{__Images/05/BW_N(20-200)_-0D_to_-4D/bw_N_20-200_Camera-2,00D(comparison).png} &
	\includegraphics[width=0.185\textwidth]{__Images/05/BW_N(20-200)_-0D_to_-4D/bw_N_20-200_Camera-3,00D(comparison).png} &
	\includegraphics[width=0.185\textwidth]{__Images/05/BW_N(20-200)_-0D_to_-4D/bw_N_20-200_Camera-4,00D(comparison).png} \\
	
	\begin{sideways} \parbox[b]{20mm} {Difference} \end{sideways} &
	\includegraphics[width=0.185\textwidth]{__Images/05/BW_N(20-200)_-0D_to_-4D/bw_N_20-200_Camera-0,00D(diff).png} &
	\includegraphics[width=0.185\textwidth]{__Images/05/BW_N(20-200)_-0D_to_-4D/bw_N_20-200_Camera-1,00D(diff).png} &
	\includegraphics[width=0.185\textwidth]{__Images/05/BW_N(20-200)_-0D_to_-4D/bw_N_20-200_Camera-2,00D(diff).png} &
	\includegraphics[width=0.185\textwidth]{__Images/05/BW_N(20-200)_-0D_to_-4D/bw_N_20-200_Camera-3,00D(diff).png} &
	\includegraphics[width=0.185\textwidth]{__Images/05/BW_N(20-200)_-0D_to_-4D/bw_N_20-200_Camera-4,00D(diff).png} \\

	\end{tabular}
	
	\caption[Comparisons of our simulated results against ground truth obtained with a hyperopic camera]{Comparisons of our simulated results against ground truth obtained with a hyperopic camera. These images correspond to a Snellen ratio of 20/200. (top row) Images captured using the DSLR camera with extra lenses varying from 0.0 to 4.0 diopters. (second row) Our simulated results. (third row) SSIM metric results. (fourth row) AD metric.}
	\label{fig:comparison_hyperopic_bw}
\end{figure*}
	% fig:comparison_myopic_bw

%\startsquarepar
Tables~\ref{table:metrics_myopic_wb} and \ref{table:metrics_myopic_bw} show the numerical results of the SSIM and PSNR metrics for the results presented in Figures~\ref{fig:comparison_myopic_wb} and \ref{fig:comparison_myopic_bw}, respectively. Each row represents the value of a specific metric (\ie, SSIM or PSNR) when comparing an image captured by the DSLR camera with the one obtained using our simulation.  
The values of the SSIM metric range from 0.0 (poor similarity) to 1.0 (high similarity). In these tables, one can see that all values are very close to 1.0, indicating that our simulations indeed produces results that are structurally very  similar the ground truth. The PSNR values also indicate that our simulations also produce results very similar to the ground truth. Note that PSNR values of 34.0 decibels and above indicate that two images are essentially indistinguishable from each other.  

%It is possible to visualize the results of this metric as an image, which represents the local (\ie, pixel-by-pixel) SSIM index. Such images are illustrated in Figures~\ref{fig:comparison_myopic_wb}(third row) and \ref{fig:comparison_myopic_bw}(third row). We've obtained the PSNR by calculating the mean squared error (MSE) and then dividing the maximum signal value that exists in our optical ground truth image by the MSE. This metric is given in decibel and relies strictly on numeric comparison between images. When two images are identical, the PSNR tends to infinity. Excellent values range from 30 to 50 dB, where higher is better.
%\stopsquarepar

% Myopic 		- Tables
% !TEX root = ./../../../_Thesis.tex

\begin{table}[!b]
	\centering
	\label{table:metrics_hyperopic_bw}
	\caption[SSIM and PSNR table of hyperopic perception (Figure~\ref{fig:comparison_hyperopic_bw})]{SSIM and PSNR table of hyperopic perception (Figure~\ref{fig:comparison_hyperopic_bw})}

	\begin{tabular}{cccccc}
	{\bf }                          & {\bf -0.00 D} & {\bf -1.00 D} & {\bf -2.00 D} & {\bf -3.00 D} & {\bf -4.00 D} \\ \hline
	\multicolumn{1}{c|}{{\bf SSIM}} & 0.9833  		& 0.9157  		& 0.9196  		& 0.9213  		& 0.9165  		\\
	\multicolumn{1}{c|}{{\bf PSNR}} & 34.6438 		& 35.0465 		& 36.6722 		& 33.3992 		& 30.9853
	\end{tabular}
\end{table}	% table:metrics_myopic_wb
% !TEX root = ./../../../_Thesis.tex

\begin{table}[!b]
	\centering
	\label{table:metrics_hyperopic_bw}
	\caption[SSIM and PSNR table of hyperopic perception (Figure~\ref{fig:comparison_hyperopic_bw})]{SSIM and PSNR table of hyperopic perception (Figure~\ref{fig:comparison_hyperopic_bw})}

	\begin{tabular}{cccccc}
	{\bf }                          & {\bf -0.00 D} & {\bf -1.00 D} & {\bf -2.00 D} & {\bf -3.00 D} & {\bf -4.00 D} \\ \hline
	\multicolumn{1}{c|}{{\bf SSIM}} & 0.9833  		& 0.9157  		& 0.9196  		& 0.9213  		& 0.9165  		\\
	\multicolumn{1}{c|}{{\bf PSNR}} & 34.6438 		& 35.0465 		& 36.6722 		& 33.3992 		& 30.9853
	\end{tabular}
\end{table}	% table:metrics_myopic_bw

%\noindent

Figures~\ref{fig:comparison_hyperopic_wb} and \ref{fig:comparison_hyperopic_bw} provide similar comparisons for hyperopic vision. The results in the top row were captured with a DSLR camera and extra lenses ranging from 0 to -4 diopters, in steps of -1 diopter. The second rows show our simulated results. Likewise, the third row presents the visualization of the pixel-by-pixel SSIM index when comparing the captured and simulated hyperopic results, and the bottom one presents the absolute difference of images. Tables~\ref{table:metrics_hyperopic_wb} and \ref{table:metrics_hyperopic_bw} provide the SSIM and PSNR values comparing the simulated images to the ground truth, attesting the quality of our results. 

% Hyperopic 	- Figures and Tables
% !TEX root = ./../../../_Thesis.tex

\begin{figure*}[!t]
	\centering

	\begin{tabular}{@{}r@{ } c@{ } c@{ } c@{ } c@{ } c }
	&
	\small{-0.00 D} &
	\small{-1.00 D} &
	\small{-2.00 D} &
	\small{-3.00 D} &
	\small{-4.00 D} & \\

	\begin{sideways} \parbox[b]{20mm} {Camera} \end{sideways} &
	\includegraphics[width=0.185\textwidth]{__Images/05/BW_N(20-200)_-0D_to_-4D/bw_N_20-200_Camera-0,00D(lens).png} &
	\includegraphics[width=0.185\textwidth]{__Images/05/BW_N(20-200)_-0D_to_-4D/bw_N_20-200_Camera-1,00D(lens).png} &
	\includegraphics[width=0.185\textwidth]{__Images/05/BW_N(20-200)_-0D_to_-4D/bw_N_20-200_Camera-2,00D(lens).png} &
	\includegraphics[width=0.185\textwidth]{__Images/05/BW_N(20-200)_-0D_to_-4D/bw_N_20-200_Camera-3,00D(lens).png} &
	\includegraphics[width=0.185\textwidth]{__Images/05/BW_N(20-200)_-0D_to_-4D/bw_N_20-200_Camera-4,00D(lens).png} \\

	\begin{sideways} \parbox[b]{20mm} {Simulation} \end{sideways} &
	\includegraphics[width=0.185\textwidth]{__Images/05/BW_N(20-200)_-0D_to_-4D/bw_N_20-200_Camera-0,00D(simulated).png} &
	\includegraphics[width=0.185\textwidth]{__Images/05/BW_N(20-200)_-0D_to_-4D/bw_N_20-200_Camera-1,00D(simulated).png} &
	\includegraphics[width=0.185\textwidth]{__Images/05/BW_N(20-200)_-0D_to_-4D/bw_N_20-200_Camera-2,00D(simulated).png} &
	\includegraphics[width=0.185\textwidth]{__Images/05/BW_N(20-200)_-0D_to_-4D/bw_N_20-200_Camera-3,00D(simulated).png} &
	\includegraphics[width=0.185\textwidth]{__Images/05/BW_N(20-200)_-0D_to_-4D/bw_N_20-200_Camera-4,00D(simulated).png} \\

	\begin{sideways} \parbox[b]{20mm} {Local~SSIM} \end{sideways} &
	\includegraphics[width=0.185\textwidth]{__Images/05/BW_N(20-200)_-0D_to_-4D/bw_N_20-200_Camera-0,00D(comparison).png} &
	\includegraphics[width=0.185\textwidth]{__Images/05/BW_N(20-200)_-0D_to_-4D/bw_N_20-200_Camera-1,00D(comparison).png} &
	\includegraphics[width=0.185\textwidth]{__Images/05/BW_N(20-200)_-0D_to_-4D/bw_N_20-200_Camera-2,00D(comparison).png} &
	\includegraphics[width=0.185\textwidth]{__Images/05/BW_N(20-200)_-0D_to_-4D/bw_N_20-200_Camera-3,00D(comparison).png} &
	\includegraphics[width=0.185\textwidth]{__Images/05/BW_N(20-200)_-0D_to_-4D/bw_N_20-200_Camera-4,00D(comparison).png} \\
	
	\begin{sideways} \parbox[b]{20mm} {Difference} \end{sideways} &
	\includegraphics[width=0.185\textwidth]{__Images/05/BW_N(20-200)_-0D_to_-4D/bw_N_20-200_Camera-0,00D(diff).png} &
	\includegraphics[width=0.185\textwidth]{__Images/05/BW_N(20-200)_-0D_to_-4D/bw_N_20-200_Camera-1,00D(diff).png} &
	\includegraphics[width=0.185\textwidth]{__Images/05/BW_N(20-200)_-0D_to_-4D/bw_N_20-200_Camera-2,00D(diff).png} &
	\includegraphics[width=0.185\textwidth]{__Images/05/BW_N(20-200)_-0D_to_-4D/bw_N_20-200_Camera-3,00D(diff).png} &
	\includegraphics[width=0.185\textwidth]{__Images/05/BW_N(20-200)_-0D_to_-4D/bw_N_20-200_Camera-4,00D(diff).png} \\

	\end{tabular}
	
	\caption[Comparisons of our simulated results against ground truth obtained with a hyperopic camera]{Comparisons of our simulated results against ground truth obtained with a hyperopic camera. These images correspond to a Snellen ratio of 20/200. (top row) Images captured using the DSLR camera with extra lenses varying from 0.0 to 4.0 diopters. (second row) Our simulated results. (third row) SSIM metric results. (fourth row) AD metric.}
	\label{fig:comparison_hyperopic_bw}
\end{figure*}
	% fig:comparison_hyperopic_wb
% !TEX root = ./../../../_Thesis.tex

\begin{table}[!b]
	\centering
	\label{table:metrics_hyperopic_bw}
	\caption[SSIM and PSNR table of hyperopic perception (Figure~\ref{fig:comparison_hyperopic_bw})]{SSIM and PSNR table of hyperopic perception (Figure~\ref{fig:comparison_hyperopic_bw})}

	\begin{tabular}{cccccc}
	{\bf }                          & {\bf -0.00 D} & {\bf -1.00 D} & {\bf -2.00 D} & {\bf -3.00 D} & {\bf -4.00 D} \\ \hline
	\multicolumn{1}{c|}{{\bf SSIM}} & 0.9833  		& 0.9157  		& 0.9196  		& 0.9213  		& 0.9165  		\\
	\multicolumn{1}{c|}{{\bf PSNR}} & 34.6438 		& 35.0465 		& 36.6722 		& 33.3992 		& 30.9853
	\end{tabular}
\end{table}	% table:metrics_hyperopic_wb
% !TEX root = ./../../../_Thesis.tex

\begin{table}[!b]
	\centering
	\label{table:metrics_hyperopic_bw}
	\caption[SSIM and PSNR table of hyperopic perception (Figure~\ref{fig:comparison_hyperopic_bw})]{SSIM and PSNR table of hyperopic perception (Figure~\ref{fig:comparison_hyperopic_bw})}

	\begin{tabular}{cccccc}
	{\bf }                          & {\bf -0.00 D} & {\bf -1.00 D} & {\bf -2.00 D} & {\bf -3.00 D} & {\bf -4.00 D} \\ \hline
	\multicolumn{1}{c|}{{\bf SSIM}} & 0.9833  		& 0.9157  		& 0.9196  		& 0.9213  		& 0.9165  		\\
	\multicolumn{1}{c|}{{\bf PSNR}} & 34.6438 		& 35.0465 		& 36.6722 		& 33.3992 		& 30.9853
	\end{tabular}
\end{table}	% table:metrics_hyperopic_bw
% !TEX root = ./../../../_Thesis.tex

\begin{figure*}[!t]
	\centering

	\begin{tabular}{@{}r@{ } c@{ } c@{ } c@{ } c@{ } c }
	&
	\small{-0.00 D} &
	\small{-1.00 D} &
	\small{-2.00 D} &
	\small{-3.00 D} &
	\small{-4.00 D} & \\

	\begin{sideways} \parbox[b]{20mm} {Camera} \end{sideways} &
	\includegraphics[width=0.185\textwidth]{__Images/05/BW_N(20-200)_-0D_to_-4D/bw_N_20-200_Camera-0,00D(lens).png} &
	\includegraphics[width=0.185\textwidth]{__Images/05/BW_N(20-200)_-0D_to_-4D/bw_N_20-200_Camera-1,00D(lens).png} &
	\includegraphics[width=0.185\textwidth]{__Images/05/BW_N(20-200)_-0D_to_-4D/bw_N_20-200_Camera-2,00D(lens).png} &
	\includegraphics[width=0.185\textwidth]{__Images/05/BW_N(20-200)_-0D_to_-4D/bw_N_20-200_Camera-3,00D(lens).png} &
	\includegraphics[width=0.185\textwidth]{__Images/05/BW_N(20-200)_-0D_to_-4D/bw_N_20-200_Camera-4,00D(lens).png} \\

	\begin{sideways} \parbox[b]{20mm} {Simulation} \end{sideways} &
	\includegraphics[width=0.185\textwidth]{__Images/05/BW_N(20-200)_-0D_to_-4D/bw_N_20-200_Camera-0,00D(simulated).png} &
	\includegraphics[width=0.185\textwidth]{__Images/05/BW_N(20-200)_-0D_to_-4D/bw_N_20-200_Camera-1,00D(simulated).png} &
	\includegraphics[width=0.185\textwidth]{__Images/05/BW_N(20-200)_-0D_to_-4D/bw_N_20-200_Camera-2,00D(simulated).png} &
	\includegraphics[width=0.185\textwidth]{__Images/05/BW_N(20-200)_-0D_to_-4D/bw_N_20-200_Camera-3,00D(simulated).png} &
	\includegraphics[width=0.185\textwidth]{__Images/05/BW_N(20-200)_-0D_to_-4D/bw_N_20-200_Camera-4,00D(simulated).png} \\

	\begin{sideways} \parbox[b]{20mm} {Local~SSIM} \end{sideways} &
	\includegraphics[width=0.185\textwidth]{__Images/05/BW_N(20-200)_-0D_to_-4D/bw_N_20-200_Camera-0,00D(comparison).png} &
	\includegraphics[width=0.185\textwidth]{__Images/05/BW_N(20-200)_-0D_to_-4D/bw_N_20-200_Camera-1,00D(comparison).png} &
	\includegraphics[width=0.185\textwidth]{__Images/05/BW_N(20-200)_-0D_to_-4D/bw_N_20-200_Camera-2,00D(comparison).png} &
	\includegraphics[width=0.185\textwidth]{__Images/05/BW_N(20-200)_-0D_to_-4D/bw_N_20-200_Camera-3,00D(comparison).png} &
	\includegraphics[width=0.185\textwidth]{__Images/05/BW_N(20-200)_-0D_to_-4D/bw_N_20-200_Camera-4,00D(comparison).png} \\
	
	\begin{sideways} \parbox[b]{20mm} {Difference} \end{sideways} &
	\includegraphics[width=0.185\textwidth]{__Images/05/BW_N(20-200)_-0D_to_-4D/bw_N_20-200_Camera-0,00D(diff).png} &
	\includegraphics[width=0.185\textwidth]{__Images/05/BW_N(20-200)_-0D_to_-4D/bw_N_20-200_Camera-1,00D(diff).png} &
	\includegraphics[width=0.185\textwidth]{__Images/05/BW_N(20-200)_-0D_to_-4D/bw_N_20-200_Camera-2,00D(diff).png} &
	\includegraphics[width=0.185\textwidth]{__Images/05/BW_N(20-200)_-0D_to_-4D/bw_N_20-200_Camera-3,00D(diff).png} &
	\includegraphics[width=0.185\textwidth]{__Images/05/BW_N(20-200)_-0D_to_-4D/bw_N_20-200_Camera-4,00D(diff).png} \\

	\end{tabular}
	
	\caption[Comparisons of our simulated results against ground truth obtained with a hyperopic camera]{Comparisons of our simulated results against ground truth obtained with a hyperopic camera. These images correspond to a Snellen ratio of 20/200. (top row) Images captured using the DSLR camera with extra lenses varying from 0.0 to 4.0 diopters. (second row) Our simulated results. (third row) SSIM metric results. (fourth row) AD metric.}
	\label{fig:comparison_hyperopic_bw}
\end{figure*}
	% fig:comparison_hyperopic_bw

Besides simulating the effects of defocus (\ie, myopia and hyperopia), we have also compared the results of our simulation for astigmatic vision.
%We've also explored simulations of more effects rather than defocus. 
This is illustrated in Figures~\ref{fig:comparison_astig-2@90_wb} and \ref{fig:comparison_astig-2@180_wb}.
% shows the simulation of an astigmatic perception. % cylindrical lens with -2 diopters at the vertical meridian
The Sloan letters in Figure~\ref{fig:comparison_astig-2@90_wb} were captured by the DSLR camera with an additional cylindrical lens with -2.0 diopters, rotated in order to simulate astigmatism in the horizontal meridian ($\phi = 90^{\circ}$). Similarly, Figure~\ref{fig:comparison_astig-2@180_wb} shows the real and simulated astigmatism in the vertical meridian ($\phi = 180^{\circ}$). 
Figures~\ref{fig:comparison_astig+2@90_bw} and \ref{fig:comparison_astig+2@180_bw} show the captured and simulated results for a cylindrical lens with +2.0 diopters. 
Tables~\ref{table:metrics_astig-2} and \ref{table:metrics_astig+2} show the results of the SSIM and PSNR metrics for these astigmatic results. Again, the SSIM indices are close 1.0 and the PSNR is close to or above 34.00 decibels.
%Both astigmatic simulations outcomes were evaluated using the SSIM and PSNR metrics, and the resultant values when comparing the simulated images to the reference ones are shown in .

Note that for the astigmatic results, part of the differences visible in 
the astigmatic local SSIM index visualizations (Figures~\ref{fig:comparison_astig-2@90_wb}, \ref{fig:comparison_astig-2@180_wb}, \ref{fig:comparison_astig+2@90_bw} and \ref{fig:comparison_astig+2@180_bw}) is due to the difficulty of a precise manual alignment of the astigmatic axes to the ones used in our simulations. Any deviation from the simulated angles affects the results of the quality metric.

% !TEX root = ./../../../_Thesis.tex

\begin{figure*}[!t]
	\centering

	\begin{tabular}{@{}r@{ } c@{ } c@{ } c@{ } c@{ } c }
	&
	\small{-0.00 D} &
	\small{-1.00 D} &
	\small{-2.00 D} &
	\small{-3.00 D} &
	\small{-4.00 D} & \\

	\begin{sideways} \parbox[b]{20mm} {Camera} \end{sideways} &
	\includegraphics[width=0.185\textwidth]{__Images/05/BW_N(20-200)_-0D_to_-4D/bw_N_20-200_Camera-0,00D(lens).png} &
	\includegraphics[width=0.185\textwidth]{__Images/05/BW_N(20-200)_-0D_to_-4D/bw_N_20-200_Camera-1,00D(lens).png} &
	\includegraphics[width=0.185\textwidth]{__Images/05/BW_N(20-200)_-0D_to_-4D/bw_N_20-200_Camera-2,00D(lens).png} &
	\includegraphics[width=0.185\textwidth]{__Images/05/BW_N(20-200)_-0D_to_-4D/bw_N_20-200_Camera-3,00D(lens).png} &
	\includegraphics[width=0.185\textwidth]{__Images/05/BW_N(20-200)_-0D_to_-4D/bw_N_20-200_Camera-4,00D(lens).png} \\

	\begin{sideways} \parbox[b]{20mm} {Simulation} \end{sideways} &
	\includegraphics[width=0.185\textwidth]{__Images/05/BW_N(20-200)_-0D_to_-4D/bw_N_20-200_Camera-0,00D(simulated).png} &
	\includegraphics[width=0.185\textwidth]{__Images/05/BW_N(20-200)_-0D_to_-4D/bw_N_20-200_Camera-1,00D(simulated).png} &
	\includegraphics[width=0.185\textwidth]{__Images/05/BW_N(20-200)_-0D_to_-4D/bw_N_20-200_Camera-2,00D(simulated).png} &
	\includegraphics[width=0.185\textwidth]{__Images/05/BW_N(20-200)_-0D_to_-4D/bw_N_20-200_Camera-3,00D(simulated).png} &
	\includegraphics[width=0.185\textwidth]{__Images/05/BW_N(20-200)_-0D_to_-4D/bw_N_20-200_Camera-4,00D(simulated).png} \\

	\begin{sideways} \parbox[b]{20mm} {Local~SSIM} \end{sideways} &
	\includegraphics[width=0.185\textwidth]{__Images/05/BW_N(20-200)_-0D_to_-4D/bw_N_20-200_Camera-0,00D(comparison).png} &
	\includegraphics[width=0.185\textwidth]{__Images/05/BW_N(20-200)_-0D_to_-4D/bw_N_20-200_Camera-1,00D(comparison).png} &
	\includegraphics[width=0.185\textwidth]{__Images/05/BW_N(20-200)_-0D_to_-4D/bw_N_20-200_Camera-2,00D(comparison).png} &
	\includegraphics[width=0.185\textwidth]{__Images/05/BW_N(20-200)_-0D_to_-4D/bw_N_20-200_Camera-3,00D(comparison).png} &
	\includegraphics[width=0.185\textwidth]{__Images/05/BW_N(20-200)_-0D_to_-4D/bw_N_20-200_Camera-4,00D(comparison).png} \\
	
	\begin{sideways} \parbox[b]{20mm} {Difference} \end{sideways} &
	\includegraphics[width=0.185\textwidth]{__Images/05/BW_N(20-200)_-0D_to_-4D/bw_N_20-200_Camera-0,00D(diff).png} &
	\includegraphics[width=0.185\textwidth]{__Images/05/BW_N(20-200)_-0D_to_-4D/bw_N_20-200_Camera-1,00D(diff).png} &
	\includegraphics[width=0.185\textwidth]{__Images/05/BW_N(20-200)_-0D_to_-4D/bw_N_20-200_Camera-2,00D(diff).png} &
	\includegraphics[width=0.185\textwidth]{__Images/05/BW_N(20-200)_-0D_to_-4D/bw_N_20-200_Camera-3,00D(diff).png} &
	\includegraphics[width=0.185\textwidth]{__Images/05/BW_N(20-200)_-0D_to_-4D/bw_N_20-200_Camera-4,00D(diff).png} \\

	\end{tabular}
	
	\caption[Comparisons of our simulated results against ground truth obtained with a hyperopic camera]{Comparisons of our simulated results against ground truth obtained with a hyperopic camera. These images correspond to a Snellen ratio of 20/200. (top row) Images captured using the DSLR camera with extra lenses varying from 0.0 to 4.0 diopters. (second row) Our simulated results. (third row) SSIM metric results. (fourth row) AD metric.}
	\label{fig:comparison_hyperopic_bw}
\end{figure*}
			% fig:comparison_astig-2@90_wb
% !TEX root = ./../../../_Thesis.tex

\begin{figure*}[!t]
	\centering

	\begin{tabular}{@{}r@{ } c@{ } c@{ } c@{ } c@{ } c }
	&
	\small{-0.00 D} &
	\small{-1.00 D} &
	\small{-2.00 D} &
	\small{-3.00 D} &
	\small{-4.00 D} & \\

	\begin{sideways} \parbox[b]{20mm} {Camera} \end{sideways} &
	\includegraphics[width=0.185\textwidth]{__Images/05/BW_N(20-200)_-0D_to_-4D/bw_N_20-200_Camera-0,00D(lens).png} &
	\includegraphics[width=0.185\textwidth]{__Images/05/BW_N(20-200)_-0D_to_-4D/bw_N_20-200_Camera-1,00D(lens).png} &
	\includegraphics[width=0.185\textwidth]{__Images/05/BW_N(20-200)_-0D_to_-4D/bw_N_20-200_Camera-2,00D(lens).png} &
	\includegraphics[width=0.185\textwidth]{__Images/05/BW_N(20-200)_-0D_to_-4D/bw_N_20-200_Camera-3,00D(lens).png} &
	\includegraphics[width=0.185\textwidth]{__Images/05/BW_N(20-200)_-0D_to_-4D/bw_N_20-200_Camera-4,00D(lens).png} \\

	\begin{sideways} \parbox[b]{20mm} {Simulation} \end{sideways} &
	\includegraphics[width=0.185\textwidth]{__Images/05/BW_N(20-200)_-0D_to_-4D/bw_N_20-200_Camera-0,00D(simulated).png} &
	\includegraphics[width=0.185\textwidth]{__Images/05/BW_N(20-200)_-0D_to_-4D/bw_N_20-200_Camera-1,00D(simulated).png} &
	\includegraphics[width=0.185\textwidth]{__Images/05/BW_N(20-200)_-0D_to_-4D/bw_N_20-200_Camera-2,00D(simulated).png} &
	\includegraphics[width=0.185\textwidth]{__Images/05/BW_N(20-200)_-0D_to_-4D/bw_N_20-200_Camera-3,00D(simulated).png} &
	\includegraphics[width=0.185\textwidth]{__Images/05/BW_N(20-200)_-0D_to_-4D/bw_N_20-200_Camera-4,00D(simulated).png} \\

	\begin{sideways} \parbox[b]{20mm} {Local~SSIM} \end{sideways} &
	\includegraphics[width=0.185\textwidth]{__Images/05/BW_N(20-200)_-0D_to_-4D/bw_N_20-200_Camera-0,00D(comparison).png} &
	\includegraphics[width=0.185\textwidth]{__Images/05/BW_N(20-200)_-0D_to_-4D/bw_N_20-200_Camera-1,00D(comparison).png} &
	\includegraphics[width=0.185\textwidth]{__Images/05/BW_N(20-200)_-0D_to_-4D/bw_N_20-200_Camera-2,00D(comparison).png} &
	\includegraphics[width=0.185\textwidth]{__Images/05/BW_N(20-200)_-0D_to_-4D/bw_N_20-200_Camera-3,00D(comparison).png} &
	\includegraphics[width=0.185\textwidth]{__Images/05/BW_N(20-200)_-0D_to_-4D/bw_N_20-200_Camera-4,00D(comparison).png} \\
	
	\begin{sideways} \parbox[b]{20mm} {Difference} \end{sideways} &
	\includegraphics[width=0.185\textwidth]{__Images/05/BW_N(20-200)_-0D_to_-4D/bw_N_20-200_Camera-0,00D(diff).png} &
	\includegraphics[width=0.185\textwidth]{__Images/05/BW_N(20-200)_-0D_to_-4D/bw_N_20-200_Camera-1,00D(diff).png} &
	\includegraphics[width=0.185\textwidth]{__Images/05/BW_N(20-200)_-0D_to_-4D/bw_N_20-200_Camera-2,00D(diff).png} &
	\includegraphics[width=0.185\textwidth]{__Images/05/BW_N(20-200)_-0D_to_-4D/bw_N_20-200_Camera-3,00D(diff).png} &
	\includegraphics[width=0.185\textwidth]{__Images/05/BW_N(20-200)_-0D_to_-4D/bw_N_20-200_Camera-4,00D(diff).png} \\

	\end{tabular}
	
	\caption[Comparisons of our simulated results against ground truth obtained with a hyperopic camera]{Comparisons of our simulated results against ground truth obtained with a hyperopic camera. These images correspond to a Snellen ratio of 20/200. (top row) Images captured using the DSLR camera with extra lenses varying from 0.0 to 4.0 diopters. (second row) Our simulated results. (third row) SSIM metric results. (fourth row) AD metric.}
	\label{fig:comparison_hyperopic_bw}
\end{figure*}
			% fig:comparison_astig-2@180_wb
%% !TEX root = ./../../../_Thesis.tex

\begin{table}[!b]
	\centering
	\label{table:metrics_hyperopic_bw}
	\caption[SSIM and PSNR table of hyperopic perception (Figure~\ref{fig:comparison_hyperopic_bw})]{SSIM and PSNR table of hyperopic perception (Figure~\ref{fig:comparison_hyperopic_bw})}

	\begin{tabular}{cccccc}
	{\bf }                          & {\bf -0.00 D} & {\bf -1.00 D} & {\bf -2.00 D} & {\bf -3.00 D} & {\bf -4.00 D} \\ \hline
	\multicolumn{1}{c|}{{\bf SSIM}} & 0.9833  		& 0.9157  		& 0.9196  		& 0.9213  		& 0.9165  		\\
	\multicolumn{1}{c|}{{\bf PSNR}} & 34.6438 		& 35.0465 		& 36.6722 		& 33.3992 		& 30.9853
	\end{tabular}
\end{table}			% table:metrics_astig-2@90_wb
%% !TEX root = ./../../../_Thesis.tex

\begin{table}[!b]
	\centering
	\label{table:metrics_hyperopic_bw}
	\caption[SSIM and PSNR table of hyperopic perception (Figure~\ref{fig:comparison_hyperopic_bw})]{SSIM and PSNR table of hyperopic perception (Figure~\ref{fig:comparison_hyperopic_bw})}

	\begin{tabular}{cccccc}
	{\bf }                          & {\bf -0.00 D} & {\bf -1.00 D} & {\bf -2.00 D} & {\bf -3.00 D} & {\bf -4.00 D} \\ \hline
	\multicolumn{1}{c|}{{\bf SSIM}} & 0.9833  		& 0.9157  		& 0.9196  		& 0.9213  		& 0.9165  		\\
	\multicolumn{1}{c|}{{\bf PSNR}} & 34.6438 		& 35.0465 		& 36.6722 		& 33.3992 		& 30.9853
	\end{tabular}
\end{table}			% table:metrics_astig-2@180_wb
% !TEX root = ./../../_Thesis.tex

\begin{table}[!b]
	\centering
	\label{table:metrics_astig-2}
	\caption[SSIM and PSNR table of the negative astigmatic perception]{SSIM and PSNR table of the negative astigmatic perception.}
	\begin{tabular}{ccccccc}
	\hline
	{\bf }                                                 & {\bf }                          & {\bf N} & {\bf C} & {\bf K} & {\bf Z} & {\bf O} \\ \hline
	\multicolumn{1}{c|}{\multirow{2}{*}{{\bf -2.00 @90}}}  & \multicolumn{1}{c|}{{\bf SSIM}} & 0.9171  & 0.9174  & 0.9193  & 0.9242  & 0.9185  \\
	\multicolumn{1}{c|}{}                                  & \multicolumn{1}{c|}{{\bf PSNR}} & 36.4269 & 36.0369 & 37.0615 & 35.9831 & 36.9738 \\ \hline
	\multicolumn{1}{c|}{\multirow{2}{*}{{\bf -2.00 @180}}} & \multicolumn{1}{c|}{{\bf SSIM}} & 0.9169  & 0.9207  & 0.9226  & 0.9220  & 0.9201  \\
	\multicolumn{1}{c|}{}                                  & \multicolumn{1}{c|}{{\bf PSNR}} & 34.6329 & 37.8456 & 37.4705 & 38.7289 & 37.8041 \\ \hline
	\end{tabular}

\end{table}					% table:metrics_astig-2


% !TEX root = ./../../../_Thesis.tex

\begin{figure*}[!t]
	\centering

	\begin{tabular}{@{}r@{ } c@{ } c@{ } c@{ } c@{ } c }
	&
	\small{-0.00 D} &
	\small{-1.00 D} &
	\small{-2.00 D} &
	\small{-3.00 D} &
	\small{-4.00 D} & \\

	\begin{sideways} \parbox[b]{20mm} {Camera} \end{sideways} &
	\includegraphics[width=0.185\textwidth]{__Images/05/BW_N(20-200)_-0D_to_-4D/bw_N_20-200_Camera-0,00D(lens).png} &
	\includegraphics[width=0.185\textwidth]{__Images/05/BW_N(20-200)_-0D_to_-4D/bw_N_20-200_Camera-1,00D(lens).png} &
	\includegraphics[width=0.185\textwidth]{__Images/05/BW_N(20-200)_-0D_to_-4D/bw_N_20-200_Camera-2,00D(lens).png} &
	\includegraphics[width=0.185\textwidth]{__Images/05/BW_N(20-200)_-0D_to_-4D/bw_N_20-200_Camera-3,00D(lens).png} &
	\includegraphics[width=0.185\textwidth]{__Images/05/BW_N(20-200)_-0D_to_-4D/bw_N_20-200_Camera-4,00D(lens).png} \\

	\begin{sideways} \parbox[b]{20mm} {Simulation} \end{sideways} &
	\includegraphics[width=0.185\textwidth]{__Images/05/BW_N(20-200)_-0D_to_-4D/bw_N_20-200_Camera-0,00D(simulated).png} &
	\includegraphics[width=0.185\textwidth]{__Images/05/BW_N(20-200)_-0D_to_-4D/bw_N_20-200_Camera-1,00D(simulated).png} &
	\includegraphics[width=0.185\textwidth]{__Images/05/BW_N(20-200)_-0D_to_-4D/bw_N_20-200_Camera-2,00D(simulated).png} &
	\includegraphics[width=0.185\textwidth]{__Images/05/BW_N(20-200)_-0D_to_-4D/bw_N_20-200_Camera-3,00D(simulated).png} &
	\includegraphics[width=0.185\textwidth]{__Images/05/BW_N(20-200)_-0D_to_-4D/bw_N_20-200_Camera-4,00D(simulated).png} \\

	\begin{sideways} \parbox[b]{20mm} {Local~SSIM} \end{sideways} &
	\includegraphics[width=0.185\textwidth]{__Images/05/BW_N(20-200)_-0D_to_-4D/bw_N_20-200_Camera-0,00D(comparison).png} &
	\includegraphics[width=0.185\textwidth]{__Images/05/BW_N(20-200)_-0D_to_-4D/bw_N_20-200_Camera-1,00D(comparison).png} &
	\includegraphics[width=0.185\textwidth]{__Images/05/BW_N(20-200)_-0D_to_-4D/bw_N_20-200_Camera-2,00D(comparison).png} &
	\includegraphics[width=0.185\textwidth]{__Images/05/BW_N(20-200)_-0D_to_-4D/bw_N_20-200_Camera-3,00D(comparison).png} &
	\includegraphics[width=0.185\textwidth]{__Images/05/BW_N(20-200)_-0D_to_-4D/bw_N_20-200_Camera-4,00D(comparison).png} \\
	
	\begin{sideways} \parbox[b]{20mm} {Difference} \end{sideways} &
	\includegraphics[width=0.185\textwidth]{__Images/05/BW_N(20-200)_-0D_to_-4D/bw_N_20-200_Camera-0,00D(diff).png} &
	\includegraphics[width=0.185\textwidth]{__Images/05/BW_N(20-200)_-0D_to_-4D/bw_N_20-200_Camera-1,00D(diff).png} &
	\includegraphics[width=0.185\textwidth]{__Images/05/BW_N(20-200)_-0D_to_-4D/bw_N_20-200_Camera-2,00D(diff).png} &
	\includegraphics[width=0.185\textwidth]{__Images/05/BW_N(20-200)_-0D_to_-4D/bw_N_20-200_Camera-3,00D(diff).png} &
	\includegraphics[width=0.185\textwidth]{__Images/05/BW_N(20-200)_-0D_to_-4D/bw_N_20-200_Camera-4,00D(diff).png} \\

	\end{tabular}
	
	\caption[Comparisons of our simulated results against ground truth obtained with a hyperopic camera]{Comparisons of our simulated results against ground truth obtained with a hyperopic camera. These images correspond to a Snellen ratio of 20/200. (top row) Images captured using the DSLR camera with extra lenses varying from 0.0 to 4.0 diopters. (second row) Our simulated results. (third row) SSIM metric results. (fourth row) AD metric.}
	\label{fig:comparison_hyperopic_bw}
\end{figure*}
			% fig:comparison_astig+2@90_bw
% !TEX root = ./../../../_Thesis.tex

\begin{figure*}[!t]
	\centering

	\begin{tabular}{@{}r@{ } c@{ } c@{ } c@{ } c@{ } c }
	&
	\small{-0.00 D} &
	\small{-1.00 D} &
	\small{-2.00 D} &
	\small{-3.00 D} &
	\small{-4.00 D} & \\

	\begin{sideways} \parbox[b]{20mm} {Camera} \end{sideways} &
	\includegraphics[width=0.185\textwidth]{__Images/05/BW_N(20-200)_-0D_to_-4D/bw_N_20-200_Camera-0,00D(lens).png} &
	\includegraphics[width=0.185\textwidth]{__Images/05/BW_N(20-200)_-0D_to_-4D/bw_N_20-200_Camera-1,00D(lens).png} &
	\includegraphics[width=0.185\textwidth]{__Images/05/BW_N(20-200)_-0D_to_-4D/bw_N_20-200_Camera-2,00D(lens).png} &
	\includegraphics[width=0.185\textwidth]{__Images/05/BW_N(20-200)_-0D_to_-4D/bw_N_20-200_Camera-3,00D(lens).png} &
	\includegraphics[width=0.185\textwidth]{__Images/05/BW_N(20-200)_-0D_to_-4D/bw_N_20-200_Camera-4,00D(lens).png} \\

	\begin{sideways} \parbox[b]{20mm} {Simulation} \end{sideways} &
	\includegraphics[width=0.185\textwidth]{__Images/05/BW_N(20-200)_-0D_to_-4D/bw_N_20-200_Camera-0,00D(simulated).png} &
	\includegraphics[width=0.185\textwidth]{__Images/05/BW_N(20-200)_-0D_to_-4D/bw_N_20-200_Camera-1,00D(simulated).png} &
	\includegraphics[width=0.185\textwidth]{__Images/05/BW_N(20-200)_-0D_to_-4D/bw_N_20-200_Camera-2,00D(simulated).png} &
	\includegraphics[width=0.185\textwidth]{__Images/05/BW_N(20-200)_-0D_to_-4D/bw_N_20-200_Camera-3,00D(simulated).png} &
	\includegraphics[width=0.185\textwidth]{__Images/05/BW_N(20-200)_-0D_to_-4D/bw_N_20-200_Camera-4,00D(simulated).png} \\

	\begin{sideways} \parbox[b]{20mm} {Local~SSIM} \end{sideways} &
	\includegraphics[width=0.185\textwidth]{__Images/05/BW_N(20-200)_-0D_to_-4D/bw_N_20-200_Camera-0,00D(comparison).png} &
	\includegraphics[width=0.185\textwidth]{__Images/05/BW_N(20-200)_-0D_to_-4D/bw_N_20-200_Camera-1,00D(comparison).png} &
	\includegraphics[width=0.185\textwidth]{__Images/05/BW_N(20-200)_-0D_to_-4D/bw_N_20-200_Camera-2,00D(comparison).png} &
	\includegraphics[width=0.185\textwidth]{__Images/05/BW_N(20-200)_-0D_to_-4D/bw_N_20-200_Camera-3,00D(comparison).png} &
	\includegraphics[width=0.185\textwidth]{__Images/05/BW_N(20-200)_-0D_to_-4D/bw_N_20-200_Camera-4,00D(comparison).png} \\
	
	\begin{sideways} \parbox[b]{20mm} {Difference} \end{sideways} &
	\includegraphics[width=0.185\textwidth]{__Images/05/BW_N(20-200)_-0D_to_-4D/bw_N_20-200_Camera-0,00D(diff).png} &
	\includegraphics[width=0.185\textwidth]{__Images/05/BW_N(20-200)_-0D_to_-4D/bw_N_20-200_Camera-1,00D(diff).png} &
	\includegraphics[width=0.185\textwidth]{__Images/05/BW_N(20-200)_-0D_to_-4D/bw_N_20-200_Camera-2,00D(diff).png} &
	\includegraphics[width=0.185\textwidth]{__Images/05/BW_N(20-200)_-0D_to_-4D/bw_N_20-200_Camera-3,00D(diff).png} &
	\includegraphics[width=0.185\textwidth]{__Images/05/BW_N(20-200)_-0D_to_-4D/bw_N_20-200_Camera-4,00D(diff).png} \\

	\end{tabular}
	
	\caption[Comparisons of our simulated results against ground truth obtained with a hyperopic camera]{Comparisons of our simulated results against ground truth obtained with a hyperopic camera. These images correspond to a Snellen ratio of 20/200. (top row) Images captured using the DSLR camera with extra lenses varying from 0.0 to 4.0 diopters. (second row) Our simulated results. (third row) SSIM metric results. (fourth row) AD metric.}
	\label{fig:comparison_hyperopic_bw}
\end{figure*}
			% fig:comparison_astig+2@180_bw
%% !TEX root = ./../../../_Thesis.tex

\begin{table}[!b]
	\centering
	\label{table:metrics_hyperopic_bw}
	\caption[SSIM and PSNR table of hyperopic perception (Figure~\ref{fig:comparison_hyperopic_bw})]{SSIM and PSNR table of hyperopic perception (Figure~\ref{fig:comparison_hyperopic_bw})}

	\begin{tabular}{cccccc}
	{\bf }                          & {\bf -0.00 D} & {\bf -1.00 D} & {\bf -2.00 D} & {\bf -3.00 D} & {\bf -4.00 D} \\ \hline
	\multicolumn{1}{c|}{{\bf SSIM}} & 0.9833  		& 0.9157  		& 0.9196  		& 0.9213  		& 0.9165  		\\
	\multicolumn{1}{c|}{{\bf PSNR}} & 34.6438 		& 35.0465 		& 36.6722 		& 33.3992 		& 30.9853
	\end{tabular}
\end{table}			% table:metrics_astig+2@90_bw
%% !TEX root = ./../../../_Thesis.tex

\begin{table}[!b]
	\centering
	\label{table:metrics_hyperopic_bw}
	\caption[SSIM and PSNR table of hyperopic perception (Figure~\ref{fig:comparison_hyperopic_bw})]{SSIM and PSNR table of hyperopic perception (Figure~\ref{fig:comparison_hyperopic_bw})}

	\begin{tabular}{cccccc}
	{\bf }                          & {\bf -0.00 D} & {\bf -1.00 D} & {\bf -2.00 D} & {\bf -3.00 D} & {\bf -4.00 D} \\ \hline
	\multicolumn{1}{c|}{{\bf SSIM}} & 0.9833  		& 0.9157  		& 0.9196  		& 0.9213  		& 0.9165  		\\
	\multicolumn{1}{c|}{{\bf PSNR}} & 34.6438 		& 35.0465 		& 36.6722 		& 33.3992 		& 30.9853
	\end{tabular}
\end{table}			% table:metrics_astig+2@180_bw
% !TEX root = ./../../_Thesis.tex

\begin{table}[!b]
	\centering
	\label{table:metrics_astig+2}
	\caption[SSIM and PSNR table of the positive astigmatic perception]{SSIM and PSNR table of the positive astigmatic perception.}
	\begin{tabular}{ccccccc}
	\hline
	{\bf }                                                 & {\bf }                          & {\bf N} & {\bf C} & {\bf K} & {\bf Z} & {\bf O} \\ \hline
	\multicolumn{1}{c|}{\multirow{2}{*}{{\bf +2.00 @90}}}  & \multicolumn{1}{c|}{{\bf SSIM}} & 0.9307  & 0.9271  & 0.9277  & 0.9193  & 0.9303  \\
	\multicolumn{1}{c|}{}                                  & \multicolumn{1}{c|}{{\bf PSNR}} & 37.2835 & 35.4130 & 36.3713 & 32.7150 & 36.2564 \\ \hline
	\multicolumn{1}{c|}{\multirow{2}{*}{{\bf +2.00 @180}}} & \multicolumn{1}{c|}{{\bf SSIM}} & 0.9235  & 0.9287  & 0.9257  & 0.9308  & 0.9260  \\
	\multicolumn{1}{c|}{}                                  & \multicolumn{1}{c|}{{\bf PSNR}} & 33.3036 & 36.8478 & 35.0677 & 36.5272 & 36.2953 \\ \hline
	\end{tabular}

\end{table}					% table:metrics_astig+2


Figures~\ref{fig:myopia_nckzo_wb} and ~\ref{fig:myopia_nckzo_bw} compare the results of our simulations with the images captured by a myopic camera (additional lens of +2.0 D) when looking to several 20/200 Sloan letters from a distance of three feet. 
The first column shows synthetic Sloan letters used as input to produce the simulated results shown in the second column ({\it Synthetic Simulation}). 
The last three columns show, respectively, images captured by the DSLR camera, images captured by the DSLR camera with an additional +2.0 D lens, and the results of our simulations for +2.0408 D (see Table~\ref{table:newDioptricPower}) of myopia applied to the images shown in the column {\it Camera Capture}.
%
% Obserrvação interessante, mas requereria uma discussão mais aprofundada
%Note that as our method involves the convolution of the image with a blurring kernel (\ie, PSF), the noise of the reference optical images (camera ones) is attenuated in our simulations (simulation ones).


% !TEX root = ./../../../_Thesis.tex

\begin{figure*}[!t]
	\centering

	\begin{tabular}{@{}r@{ } c@{ } c@{ } c@{ } c@{ } c }
	&
	\small{-0.00 D} &
	\small{-1.00 D} &
	\small{-2.00 D} &
	\small{-3.00 D} &
	\small{-4.00 D} & \\

	\begin{sideways} \parbox[b]{20mm} {Camera} \end{sideways} &
	\includegraphics[width=0.185\textwidth]{__Images/05/BW_N(20-200)_-0D_to_-4D/bw_N_20-200_Camera-0,00D(lens).png} &
	\includegraphics[width=0.185\textwidth]{__Images/05/BW_N(20-200)_-0D_to_-4D/bw_N_20-200_Camera-1,00D(lens).png} &
	\includegraphics[width=0.185\textwidth]{__Images/05/BW_N(20-200)_-0D_to_-4D/bw_N_20-200_Camera-2,00D(lens).png} &
	\includegraphics[width=0.185\textwidth]{__Images/05/BW_N(20-200)_-0D_to_-4D/bw_N_20-200_Camera-3,00D(lens).png} &
	\includegraphics[width=0.185\textwidth]{__Images/05/BW_N(20-200)_-0D_to_-4D/bw_N_20-200_Camera-4,00D(lens).png} \\

	\begin{sideways} \parbox[b]{20mm} {Simulation} \end{sideways} &
	\includegraphics[width=0.185\textwidth]{__Images/05/BW_N(20-200)_-0D_to_-4D/bw_N_20-200_Camera-0,00D(simulated).png} &
	\includegraphics[width=0.185\textwidth]{__Images/05/BW_N(20-200)_-0D_to_-4D/bw_N_20-200_Camera-1,00D(simulated).png} &
	\includegraphics[width=0.185\textwidth]{__Images/05/BW_N(20-200)_-0D_to_-4D/bw_N_20-200_Camera-2,00D(simulated).png} &
	\includegraphics[width=0.185\textwidth]{__Images/05/BW_N(20-200)_-0D_to_-4D/bw_N_20-200_Camera-3,00D(simulated).png} &
	\includegraphics[width=0.185\textwidth]{__Images/05/BW_N(20-200)_-0D_to_-4D/bw_N_20-200_Camera-4,00D(simulated).png} \\

	\begin{sideways} \parbox[b]{20mm} {Local~SSIM} \end{sideways} &
	\includegraphics[width=0.185\textwidth]{__Images/05/BW_N(20-200)_-0D_to_-4D/bw_N_20-200_Camera-0,00D(comparison).png} &
	\includegraphics[width=0.185\textwidth]{__Images/05/BW_N(20-200)_-0D_to_-4D/bw_N_20-200_Camera-1,00D(comparison).png} &
	\includegraphics[width=0.185\textwidth]{__Images/05/BW_N(20-200)_-0D_to_-4D/bw_N_20-200_Camera-2,00D(comparison).png} &
	\includegraphics[width=0.185\textwidth]{__Images/05/BW_N(20-200)_-0D_to_-4D/bw_N_20-200_Camera-3,00D(comparison).png} &
	\includegraphics[width=0.185\textwidth]{__Images/05/BW_N(20-200)_-0D_to_-4D/bw_N_20-200_Camera-4,00D(comparison).png} \\
	
	\begin{sideways} \parbox[b]{20mm} {Difference} \end{sideways} &
	\includegraphics[width=0.185\textwidth]{__Images/05/BW_N(20-200)_-0D_to_-4D/bw_N_20-200_Camera-0,00D(diff).png} &
	\includegraphics[width=0.185\textwidth]{__Images/05/BW_N(20-200)_-0D_to_-4D/bw_N_20-200_Camera-1,00D(diff).png} &
	\includegraphics[width=0.185\textwidth]{__Images/05/BW_N(20-200)_-0D_to_-4D/bw_N_20-200_Camera-2,00D(diff).png} &
	\includegraphics[width=0.185\textwidth]{__Images/05/BW_N(20-200)_-0D_to_-4D/bw_N_20-200_Camera-3,00D(diff).png} &
	\includegraphics[width=0.185\textwidth]{__Images/05/BW_N(20-200)_-0D_to_-4D/bw_N_20-200_Camera-4,00D(diff).png} \\

	\end{tabular}
	
	\caption[Comparisons of our simulated results against ground truth obtained with a hyperopic camera]{Comparisons of our simulated results against ground truth obtained with a hyperopic camera. These images correspond to a Snellen ratio of 20/200. (top row) Images captured using the DSLR camera with extra lenses varying from 0.0 to 4.0 diopters. (second row) Our simulated results. (third row) SSIM metric results. (fourth row) AD metric.}
	\label{fig:comparison_hyperopic_bw}
\end{figure*}
			% fig:myopia_nckzo_wb
% !TEX root = ./../../../_Thesis.tex

\begin{figure*}[!t]
	\centering

	\begin{tabular}{@{}r@{ } c@{ } c@{ } c@{ } c@{ } c }
	&
	\small{-0.00 D} &
	\small{-1.00 D} &
	\small{-2.00 D} &
	\small{-3.00 D} &
	\small{-4.00 D} & \\

	\begin{sideways} \parbox[b]{20mm} {Camera} \end{sideways} &
	\includegraphics[width=0.185\textwidth]{__Images/05/BW_N(20-200)_-0D_to_-4D/bw_N_20-200_Camera-0,00D(lens).png} &
	\includegraphics[width=0.185\textwidth]{__Images/05/BW_N(20-200)_-0D_to_-4D/bw_N_20-200_Camera-1,00D(lens).png} &
	\includegraphics[width=0.185\textwidth]{__Images/05/BW_N(20-200)_-0D_to_-4D/bw_N_20-200_Camera-2,00D(lens).png} &
	\includegraphics[width=0.185\textwidth]{__Images/05/BW_N(20-200)_-0D_to_-4D/bw_N_20-200_Camera-3,00D(lens).png} &
	\includegraphics[width=0.185\textwidth]{__Images/05/BW_N(20-200)_-0D_to_-4D/bw_N_20-200_Camera-4,00D(lens).png} \\

	\begin{sideways} \parbox[b]{20mm} {Simulation} \end{sideways} &
	\includegraphics[width=0.185\textwidth]{__Images/05/BW_N(20-200)_-0D_to_-4D/bw_N_20-200_Camera-0,00D(simulated).png} &
	\includegraphics[width=0.185\textwidth]{__Images/05/BW_N(20-200)_-0D_to_-4D/bw_N_20-200_Camera-1,00D(simulated).png} &
	\includegraphics[width=0.185\textwidth]{__Images/05/BW_N(20-200)_-0D_to_-4D/bw_N_20-200_Camera-2,00D(simulated).png} &
	\includegraphics[width=0.185\textwidth]{__Images/05/BW_N(20-200)_-0D_to_-4D/bw_N_20-200_Camera-3,00D(simulated).png} &
	\includegraphics[width=0.185\textwidth]{__Images/05/BW_N(20-200)_-0D_to_-4D/bw_N_20-200_Camera-4,00D(simulated).png} \\

	\begin{sideways} \parbox[b]{20mm} {Local~SSIM} \end{sideways} &
	\includegraphics[width=0.185\textwidth]{__Images/05/BW_N(20-200)_-0D_to_-4D/bw_N_20-200_Camera-0,00D(comparison).png} &
	\includegraphics[width=0.185\textwidth]{__Images/05/BW_N(20-200)_-0D_to_-4D/bw_N_20-200_Camera-1,00D(comparison).png} &
	\includegraphics[width=0.185\textwidth]{__Images/05/BW_N(20-200)_-0D_to_-4D/bw_N_20-200_Camera-2,00D(comparison).png} &
	\includegraphics[width=0.185\textwidth]{__Images/05/BW_N(20-200)_-0D_to_-4D/bw_N_20-200_Camera-3,00D(comparison).png} &
	\includegraphics[width=0.185\textwidth]{__Images/05/BW_N(20-200)_-0D_to_-4D/bw_N_20-200_Camera-4,00D(comparison).png} \\
	
	\begin{sideways} \parbox[b]{20mm} {Difference} \end{sideways} &
	\includegraphics[width=0.185\textwidth]{__Images/05/BW_N(20-200)_-0D_to_-4D/bw_N_20-200_Camera-0,00D(diff).png} &
	\includegraphics[width=0.185\textwidth]{__Images/05/BW_N(20-200)_-0D_to_-4D/bw_N_20-200_Camera-1,00D(diff).png} &
	\includegraphics[width=0.185\textwidth]{__Images/05/BW_N(20-200)_-0D_to_-4D/bw_N_20-200_Camera-2,00D(diff).png} &
	\includegraphics[width=0.185\textwidth]{__Images/05/BW_N(20-200)_-0D_to_-4D/bw_N_20-200_Camera-3,00D(diff).png} &
	\includegraphics[width=0.185\textwidth]{__Images/05/BW_N(20-200)_-0D_to_-4D/bw_N_20-200_Camera-4,00D(diff).png} \\

	\end{tabular}
	
	\caption[Comparisons of our simulated results against ground truth obtained with a hyperopic camera]{Comparisons of our simulated results against ground truth obtained with a hyperopic camera. These images correspond to a Snellen ratio of 20/200. (top row) Images captured using the DSLR camera with extra lenses varying from 0.0 to 4.0 diopters. (second row) Our simulated results. (third row) SSIM metric results. (fourth row) AD metric.}
	\label{fig:comparison_hyperopic_bw}
\end{figure*}
			% fig:myopia_nckzo_bw

Our technique can be used to simulate arbitrary wavefront aberrations, given the corresponding aberration function $W_{(x,y}$ (Equation~\ref{eq:W}). Thus, even though such a validation depends on the existence of an optical ground truth, the method is not limited to what can be modeled using a DSLR camera and additional lenses.
%the ones configured in previous camera's setup. 
Columns {\it Aberrated Wavefront} and {\it Spatial PSF} in Figure \ref{fig:synthetic_sims} show the normalized aberrated wavefront and the spatial PSF associated with the simulation results shown in the last column for a given input letter. 
%Figure \ref{fig:synthetic_sims} (top row) 
Its top row shows how a combination of low-order aberrations (myopia and astigmatism) affects the perception of a Sloan letter. 
The second and third rows simulate, respectively, higher values of pure astigmatism and spherical aberration than one can capture with the lenses available in our trial lens set. Finally, the bottom row shows the results of a simulation involving only higher-order aberrations. 
%In this example, we have generated arbitrary values for some of the wavefront coefficients ($j=6$, $j=7$, $j=9$, $j=13$, $j=15$ and $j=18$) described in Section~\ref{subsec:VisualAberrations}.


% bw comparison figure
%
\begin{figure}[h]
	\centering

	% C
	\subfigure[]{
		\includegraphics[height=0.20\linewidth]{}
	}
	~
	\subfigure[]{
		\includegraphics[height=0.20\linewidth]{}
	}
	~
	\subfigure[]{
		\includegraphics[height=0.20\linewidth]{}
	}
	~
	\subfigure[]{
		\includegraphics[height=0.20\linewidth]{}
	}
	% O
	\subfigure[]{
		\includegraphics[height=0.20\linewidth]{}
	}
	~
	\subfigure[]{
		\includegraphics[height=0.20\linewidth]{}
	}
	~
	\subfigure[]{
		\includegraphics[height=0.20\linewidth]{}
	}
	~
	\subfigure[]{
		\includegraphics[height=0.20\linewidth]{}
	}
	% R
	\subfigure[]{
		\includegraphics[height=0.20\linewidth]{}
	}
	~
	\subfigure[]{
		\includegraphics[height=0.20\linewidth]{}
	}
	~
	\subfigure[]{
		\includegraphics[height=0.20\linewidth]{}
	}
	~
	\subfigure[]{
		\includegraphics[height=0.20\linewidth]{}
	}
	% S
	\subfigure[]{
		\includegraphics[height=0.20\linewidth]{}
	}
	~
	\subfigure[]{
		\includegraphics[height=0.20\linewidth]{}
	}
	~
	\subfigure[]{
		\includegraphics[height=0.20\linewidth]{}
	}
	~
	\subfigure[]{
		\includegraphics[height=0.20\linewidth]{}
	}
	\caption[Simulation of how lower and higher order visual aberrations affect the perception]{Simulation of how lower and higher order visual aberrations affect the perception. (a) Sloan C letter, V.A. of 20/200; (b) normalized wavefront with low-order aberrations --- +0.5 Sph. -2.0 Cyl. at 45; (c) spatial PSF; (d) final image after the convolution between images (a) and (c); (e) Sloan O letter, V.A. of 20/200; (f) normalized wavefront with low-order aberrations --- -4.7 Cyl. at 135; (g) spatial PSF; (h) final image after the convolution between images (e) and (g); (i) Sloan R letter, V.A. of 20/200; (j) normalized wavefront with low-order aberrations --- +6 Sph; (k) spatial PSF; (l) final image after the convolution between images (i) and (k); (m) Sloan S letter, V.A. of 20/200; (n) normalized wavefront with high-order aberrations --- $Z_{c=6}=0.2$, $Z_{c=7}=0.2$, $Z_{c=9}=0.1$, $Z_{c=13}=0.2$, $Z_{c=15}=0.4$ and $Z_{c=18}=0.3$; (o) spatial PSF; and (p) final image after the convolution between images (m) and (o).}
	\label{fig:synthetic_sims}
\end{figure}

\clearpage
% !TEX root = ./../../../_Thesis.tex

\begin{figure*}[h]
	\centering

	\begin{tabular}{@{}r@{ } c@{ } c@{ } c@{ } c@{ } c }
	&
	\small{Input Letter} &
	\small{Aberrated Wavefront} &
	\small{Spatial PSF} &
	\small{Simulation} & \\ \\

	\begin{sideways} \parbox[b]{25mm} {} \end{sideways} &
	\includegraphics[width= 0.22\textwidth]{__Images/05/synthetic_sims/C_20-200@4x.png} &
	\includegraphics[height=0.22\textwidth]{__Images/05/synthetic_sims/Wavefront_0,5D,-2@45.png} &
	\includegraphics[width= 0.22\textwidth]{__Images/05/synthetic_sims/PSF_0,5D,-2@45.png} &
	\includegraphics[width= 0.22\textwidth]{__Images/05/synthetic_sims/C_20-200_f50_simulated(0,5D,-2@45).png} 			\\ \\

	\begin{sideways} \parbox[b]{25mm} {} \end{sideways} &
	\includegraphics[width= 0.22\textwidth]{__Images/05/synthetic_sims/O_20-200@4x.png} &
	\includegraphics[height=0.22\textwidth]{__Images/05/synthetic_sims/Wavefront_0D,-4,7@135.png} &
	\includegraphics[width= 0.22\textwidth]{__Images/05/synthetic_sims/PSF_0D,-4,7@135.png} &
	\includegraphics[width= 0.22\textwidth]{__Images/05/synthetic_sims/O_20-200_f50_simulated(0D,-4,7@135).png}			\\ \\

	\begin{sideways} \parbox[b]{25mm} {} \end{sideways} &
	\includegraphics[width= 0.22\textwidth]{__Images/05/synthetic_sims/R_20-200@4x.png} &
	\includegraphics[height=0.22\textwidth]{__Images/05/synthetic_sims/Wavefront_6D,0@0.png} &
	\includegraphics[width= 0.22\textwidth]{__Images/05/synthetic_sims/PSF_6D,0@0.png} &
	\includegraphics[width= 0.22\textwidth]{__Images/05/synthetic_sims/R_20-200_f50_simulated(6D,0@0).png}				\\ \\
	
	\begin{sideways} \parbox[b]{25mm} {} \end{sideways} &
	\includegraphics[width= 0.22\textwidth]{__Images/05/synthetic_sims/S_20-200@4x.png} &
	\includegraphics[height=0.22\textwidth]{__Images/05/synthetic_sims/Wavefront_highorder.png} &
	\includegraphics[width= 0.22\textwidth]{__Images/05/synthetic_sims/PSF_highorder.png} &
	\includegraphics[width= 0.22\textwidth]{__Images/05/synthetic_sims/S_20-200_f50_simulated(high-order).png}			\\

	\end{tabular}
	
	\caption[Simulations with arbitrary wavefronts]{Simulations with arbitrary wavefronts. The input letter images correspond to a Snellen ratio of 20/200. (second column) Normalized aberrated wavefront. (third column) The spatial PSF. (fourth column) Our simulation results given the images shown in column \emph{Input Letter}. The top row shows how a combination of low-order aberrations (+0.5 Sph. -2.0 Cyl. at 45$^{\circ}$) affects the perception of a Sloan letter. The second and third rows simulate, respectively, higher values of pure astigmatism and spherical aberration (-4.7 Cyl. at 135$^{\circ}$ and +6 Sph.) than one can capture with the lenses available in our trial lens set. The bottom row shows the results of a simulation involving only higher-order aberrations ($Z^{-3}_{3}=0.2$, $Z^{-1}_{3}=0.2$, $Z^{3}_{3}=0.1$, $Z^{2}_{4}=0.2$, $Z^{-5}_{5}=0.4$, $Z^{1}_{5}=0.3$).}
	\label{fig:synthetic_sims}
\end{figure*}

\clearpage
	% fig:synthetic_sims

% simulation 1:		+0.5 Sph.	-2.0 Cyl.	45		WB
% simulation 2:		+0.0 Sph.	-4.0 Cyl	135		WB
% simulation 3:		+6.0 Sph	 0.0 Cyl	0		WB
% simulation 4:		high-order aberration			WB
		%coeff(13) = 0.2;
		%coeff(9)  = 0.1;
		%coeff(18) = 0.3;
		%coeff(7) =  0.2;
		%coeff(15) = 0.4;
		%coeff(6)  = 0.2;
		
\section{Summary}
%
This chapter described the visual simulation technique that represents one of the central aspects of this thesis. For this, we have presented all the  involved mathematical and optical concepts. The chapter also presented a validation of our simulation technique by comparing its results with images captured by a camera instrumented with additional lenses to induce myopia, hyperopia, and astigmatism. The results of the SSIM and PSNR metric confirm that the results produced by our simulations are structurally and perceptually similar to the ground truths captured by the camera.  

	
% ---  ABSOLUTE THRESHOLD MATERIAL

